\DocumentMetadata{
  pdfversion=2.0,
  pdfstandard=ua-2,
  lang=en-US,
  tagging=on,
  tagging-setup={math/setup=mathml-SE,tabsorder=structure}
}
\documentclass[11pt,letterpaper]{article}
\usepackage[margin=0.68in,includefoot]{geometry}
\usepackage{newcomputermodern}
\usepackage{amsmath,amscd,mathtools}
\usepackage{tikz,adjustbox}
\providecommand{\square}{\mdlgwhtsquare}
\usepackage[hidelinks]{hyperref}
\hypersetup{
  pdftitle={Numerical Linear Algebra First-Year Exam, January 2016},
  pdfauthor={University of Florida Department of Mathematics},
  pdfsubject={Graduate mathematics examination},
  pdfdisplaydoctitle=true,
  bookmarksopen=true
}
\setcounter{secnumdepth}{0}
\setlength{\parindent}{0pt}
\setlength{\parskip}{0.28em}
\setlength{\emergencystretch}{3em}
\setlength{\leftmargini}{2.15em}
\setlength{\leftmarginii}{2.65em}
\setlength{\leftmarginiii}{3em}
\pagestyle{empty}
\raggedbottom
\allowdisplaybreaks
\NewTaggingSocketPlug{block/list/label}{examproblem}{%
  \tagstructbegin{tag=Lbl}\tagstructend%
  \tagstructbegin{tag=\LBody}%
  \tagstructbegin{tag=\ProblemHeadingTag,title-o={\ProblemHeadingText}}%
  \tagmcbegin{tag=\ProblemHeadingTag,actualtext={\ProblemHeadingText}}%
  #2%
  \tagmcend%
  \tagstructend%
}
\newenvironment{examproblems}[2]{%
  \def\ProblemHeadingTag{#2}%
  \AssignTaggingSocketPlug{block/list/label}{examproblem}%
  \begin{enumerate}%
  \setcounter{enumi}{#1}%
  \renewcommand{\labelenumi}{\textbf{\theenumi.}}%
  \setlength{\topsep}{0.35em}%
  \setlength{\partopsep}{0pt}%
  \setlength{\itemsep}{0.48em}%
  \setlength{\parsep}{0.18em}%
}{\end{enumerate}}
\newenvironment{examsubparts}{%
  \AssignTaggingSocketPlug{block/list/label}{default}%
  \begin{enumerate}%
  \setlength{\topsep}{0.18em}%
  \setlength{\partopsep}{0pt}%
  \setlength{\itemsep}{0.16em}%
  \setlength{\parsep}{0.1em}%
}{\end{enumerate}}
\newenvironment{examinstructions}{%
  \par\begingroup\itshape
}{%
  \par\endgroup\vspace{0.18em}
}
\makeatletter
\renewcommand\subsection{\@startsection{subsection}{2}{0pt}%
  {-1.25ex plus -.25ex minus -.1ex}%
  {0.55ex plus .1ex}%
  {\normalfont\large\bfseries}}
\renewcommand{\@maketitle}{%
  \newpage\null\vskip -1em%
  {\footnotesize\bfseries
    \href{https://www.ufl.edu/}{University of Florida}, \href{https://math.ufl.edu/}{Department of Mathematics}\par}%
  \vskip -0.5em%
  \tagstructbegin{tag=H1,title={Numerical Linear Algebra First-Year Exam, January 2016}}%
  \tagpdfparaOff%
  \tagmcbegin{tag=H1}%
    {\LARGE\bfseries\href{https://gma.math.ufl.edu/exams/numerical-linear-algebra-first-year/}{Numerical Linear Algebra First-Year Exam}, January 2016\par}%
  \tagmcend%
  \tagpdfparaOn%
  \tagstructend%
  \vskip 0.25em%
  \tagmcbegin{artifact}%
    \hrule height 1pt%
  \tagmcend%
  \vskip 1em%
}
\makeatother
\title{Numerical Linear Algebra First-Year Exam, January 2016}
\author{}
\date{}
\begin{document}
\maketitle
\thispagestyle{empty}
\begin{examinstructions}
Do 4 (four) problems.
\end{examinstructions}
\begin{examproblems}{0}{H2}
\def\ProblemHeadingText{Problem 1.}
\item
Prove that every square matrix~\(A\) has a Schur factorization.
\def\ProblemHeadingText{Problem 2.}
\item
Given a matrix \(A \in \mathbb{C}^{m \times n}\), let
\[
B=\begin{pmatrix}0&A^*\\A&0\end{pmatrix}
\quad\text{and}\quad C=A^*A.
\]
\begin{examsubparts}
\item[\textnormal{(a)}]
Show that the singular values of~\(A\) are the absolute values of eigenvalues of~\(B\).
\item[\textnormal{(b)}]
Show that the singular values of~\(A\) are the square roots of eigenvalues of~\(C\).
\item[\textnormal{(c)}]
Assume now that~\(A\) is square and invertible. Compute the two-norm condition numbers of~\(B\) and~\(C\) in terms of the two-norm condition number of~\(A\).
\item[\textnormal{(d)}]
If~\(A\) has condition number bigger than one, which has the larger condition number, \(B\) or~\(C\)?
\end{examsubparts}
\def\ProblemHeadingText{Problem 3.}
\item
This problem has two parts.
\begin{examsubparts}
\item[\textnormal{(a)}]
Compute \(\det(\lambda I+uv^*)\) when \(\lambda \in \mathbb{C}\), \(I\) is the \(m \times m\) identity matrix, and \(u,v \in \mathbb{C}^m\).
\item[\textnormal{(b)}]
Prove necessary and sufficient conditions for \(I+uv^*\) to be nonsingular and when it is, give a formula for its inverse.
\end{examsubparts}
\def\ProblemHeadingText{Problem 4.}
\item
This problem has two parts.
\begin{examsubparts}
\item[\textnormal{(a)}]
Given Cholesky decomposition of the Hermitian positive definite matrix \(A=R^*R\), prove that \(\lVert R\rVert_2=\lVert R^*\rVert_2=\lVert A\rVert_2^{1/2}\).
\item[\textnormal{(b)}]
Now assume that~\(B\) is a matrix that can be expressed as \(B=T^*T\) for some upper triangular matrix~\(T\). Show that~\(B\) is Hermitian and positive semi-definite, i.e. \(x^*Bx \geq 0\) for all~\(x\).
\end{examsubparts}
\def\ProblemHeadingText{Problem 5.}
\item
Let \(\{q_1,q_2,\ldots,q_n\}\) be an orthonormal subset of \(\mathbb{C}^m\). Show that
\[
P=\sum_{i=1}^n q_iq_i^*
\]
 is an orthogonal projector with range equal to the span of \(\{q_1,q_2,\ldots,q_n\}\).
\end{examproblems}
\end{document}
